\scanpage{55}
\begin{flushright}
atkin \mspage{15}
\end{flushright}

\reversemarginpar the random access involv\margin{\inserted{\handwritten{e}}\handwritten{/}}\deleted{\unsure{d|a|e}}d in using Euler's series to multiply (divide) \uline{$g(x$}$)$ meant that almost all the terms of \uline{$g(x$}$)$ had to be stored in core. In practice $N=43000$ was about the limit.

The main programs for each prime \uline{$q$} were written in Fortran calling on the coded subroutines described above. The only additional feature was a Fortran function $\mathrm{IJMODK}(\mathrm{K},\mathrm{C},\mathrm{I},\mathrm{J})$ which obtains the residue of $\mathrm{I}\ldotp\mathrm{J}\ (\mathrm{mod}\ \mathrm{K})$ and overwrites its calling sequence in the\normalmarginpar\margin{\handwritten{\unsure{''}}} execution run,which Dr. Don Russell ingeniously devised for me. This allows both $\mathrm{I}$ and $\mathrm{J}$ to be as large as $6.7\times10^{10}$,whereas the direct use of Fortran would limit $\mathrm{I}\ldotp\mathrm{J}$ to this size,and also be about three times as slow. It might be appropriate here to express a hope that computer design should allow for the easy extraction of an integer \uline{\unsure{remainder|remain\`{d}er}} after division,as \reversemarginpar w\margin{\inserted{\handwritten{e}}\handwritten{/}}\deleted{\unsure{d|a|e}}ll as the quotient mainly required in floating point work.

Finally it may be of interest that some of these programs when run in 1964 exposed a low-probability \handwritten{\tikz[x=1em,y=1em,baseline=-.1em]{\draw (0,-.4)--(.35,.5);}}\inserted{timing} fault in the machine's hardware. Whereas a convergent process of numerical integration(for example) may not be aff\deleted{t}\unsure{e|c}cted if a fault occurs once in a million times, arithmetic modulo \uline{$m$} depending on remainders rather than quotients is totally destroyed.
\normalmarginpar
% archive sheet number: 54
